
%%%===========================================

% ===
\section{Final Remarks}
% ===

This thesis addresses a central limitation in the way environmental problems are commonly represented in economic reasoning: the tendency to treat the environment as an external constraint, a source of externalities, or a corrective policy problem, rather than as an endogenous component of production decisions. The research does not claim that economic theory has ignored nature. Rather, it argues that the environment is often positioned outside the agent's productive calculus, which helps explain why preservation is frequently perceived as a legal obligation, compliance cost, or moral requirement instead of a rational economic strategy.

The thesis therefore proposes the concept of \emph{economic commensalism} as a theoretical framework for investigating whether environmental preservation can become compatible with private productivity and profitability. The central intuition is that ecological preservation may become individually rational when environmental variables are incorporated into production theory and when incentives are designed according to spatial, intertemporal, and strategic constraints. In this sense, the research seeks to move beyond the opposition between economic development and environmental protection by examining the conditions under which preservation may alter the structural payoff of economic agents.

The first chapter develops the theoretical foundations of this framework. It discusses how production theory, environmental economics, common-pool resource theory, and institutional approaches have treated the relationship between economic activity and nature. It also introduces a typology of economy--environment interactions, distinguishing predation, extraction, extractivism, strict preservation, and economic commensalism. This conceptual step is necessary to clarify that not every productive use of nature is necessarily predatory, but that productive use becomes problematic when it undermines the ecological conditions that sustain future production.

The second chapter translates this theoretical framework into a formal model. Its purpose is to examine whether environmental preservation can be represented as an intertemporal and strategic choice under uncertainty, bounded rationality, ecological feedback, and institutional incentives. The model is expected to identify the conditions under which preservation may become individually rational and strategically stable. In this part of the thesis, the relevant question is not whether preservation is normatively desirable, but whether it can become a rational strategy for agents operating within a properly calibrated institutional environment.

The third chapter evaluates the empirical plausibility of the proposed framework using the Brazilian state of Mato Grosso as a case study. Mato Grosso is a relevant context because it combines high agricultural productivity, deforestation pressure, ecological heterogeneity, and the presence of different environmental governance instruments. By examining the interaction between land use, productivity, deforestation, ecological conditions, and incentive mechanisms, the empirical chapter seeks to assess whether the theoretical propositions developed in the previous chapters are compatible with observed territorial dynamics.

The expected contribution of the thesis is threefold. Theoretically, it seeks to reposition environmental preservation within production theory by treating ecological systems as potential cofactors of production rather than merely as external constraints. Formally, it aims to identify the conditions under which preservation may become compatible with private rationality and strategic equilibrium. Empirically, it intends to evaluate whether this mechanism is plausible in a concrete territory where production and preservation pressures coexist.

The thesis also has potential implications for public policy and institutional design. If environmental preservation can be incorporated into the agent's productive calculus, then instruments such as carbon markets, payments for environmental services, environmental certifications, rural credit, fiscal incentives, monitoring systems, and enforcement mechanisms should be designed not only to compensate or punish agents, but to alter the expected payoff associated with preservation. This does not imply replacing regulation with market incentives. Rather, it suggests the need for mixed arrangements in which credible enforcement prevents predatory behavior while positive incentives make preservation economically attractive.

At the same time, the argument remains conditional. Economic commensalism is not presented as a universal solution to environmental degradation. Its feasibility depends on ecological characteristics, institutional capacity, monitoring technologies, property-rights clarity, enforcement credibility, market structures, and the intertemporal preferences of agents. The thesis therefore does not aim to prove that preservation will always be profitable, but to investigate whether, under specific conditions, it can become a rational, incentive-compatible, and strategically stable component of economic behavior.

In this sense, the broader purpose of the research is to contribute to a more precise understanding of the relationship between production, incentives, and ecological preservation. Rather than asking only whether economic agents should preserve the environment, the thesis asks whether institutions can be designed so that preservation becomes part of what rational agents have reason to do. If this proposition proves theoretically coherent, formally consistent, and empirically plausible, economic commensalism may offer a useful framework for rethinking environmental governance in territories where economic development and ecological resilience are deeply interdependent.

\newpage

% fim %

